\section{Obsługa modułu radiowego}
Obsługa modułu radiowego realizowana jest poprzez autorską 
bibliotekę \texttt{ook\_433mhz} realizującą przesył i~odczyt
danych z~prostych nadajników radiowych, używając metody
\textit{On--Off Keying} (OOK).
Biblioteka jest mocno inspirowana biblioteką
\texttt{RadioHead}, ze szczególnym naciskiem na zawartą w~niej 
klasę \texttt{RH\_ASK}.

Ten podrozdział jest poświęcony obsłudze nadajnika radiowego
(struktura \texttt{Transmitter} biblioteki \texttt{ook\_433mhz})
i~nie będzie zagłebiał się w~implementacje obsługi odbiornika.

\subsection{\textit{On-Off Keying}}

\subsection{Przebieg transmisji}
Transmisję danych można podzielić na 3~etapy: synchronizację,
wysłanie bajtu startu i~wysłanie ramki danych.
Sygnał synchronizujący składa się z~3~bajtów wypełnionych 
naprzemiennie bitami o~wartościach~1~i~0 i~służy do 
przygotowania odbiornika na odczyt danych.
Bajt startu (zdefiniowany w~bibliotece stałą \texttt{START\_BYTE})
przyjmuje formę liczby o~zapisie binarnym \texttt{11001100}.
Struktura ramki obejmuje pole długości, zawierające informację 
o~liczbie przesyłanych bajtów, oraz następujące po nim 
dane właściwe.

\subsection{Struktura \texttt{Transmitter}}
